\section{Analysis Plan}
\label{sec:plan}

Choosing the statistics after seeing the data invites picking the one that
gives the nicest answer. The plan was written, committed and tagged before
any campaign data; later changes are dated amendments with their reason. It
is eight decisions.

\begin{description}[leftmargin=2.5em]
\item[Unit.] The run, one fresh process. Iterations inside a run are not
independent observations (they share the process's memory layout, its
compiled code and its thread placement), so treating them as such would be
pseudo-replication~\cite{georges2007,kalibera2013}. Run times are positive
and right-skewed and the questions are about ratios, so the analysis works
on log times and reports ratios.
\item[D1, per run.] The median of its 10 measured iterations.
\item[D2, per configuration.] A configuration is (environment, workload,
input, $k$, $N$). Its central value $T$ is the geometric mean of its 30 run
values. Speedup is $T(1)/T(N)$ within the same environment.
\item[D3, intervals.] Percentile bootstrap, 10\,000 resamples, fixed seed.
Runs are resampled, never iterations, and groups that are different
processes ($N=1$ and $N$; atomic and local; native and browser) are
resampled independently, since they are not paired.
\item[D4, RQ1.] The ratio of speedups $S(\mathrm{browser})/S(\mathrm{native})$
for each browser, workload and $N$; 1 means the browser scales like native.
\item[D5, RQ2.] For each environment $e$, $r(e,k,N) = \log T_{\mathrm{atomic}}
- \log T_{\mathrm{local}}$: how much slower the shared atomics are than the
lock-free control on the same input. The estimand is the interaction
$\Delta r = r(\mathrm{browser}) - r(\mathrm{native})$ at the same $k$ and
$N$, a difference of differences. Comparing within one input is what
removes the input's own effect on both variants.
\item[D6, effect sizes.] The ratios themselves: the speedup ratio for RQ1
and $\exp(\Delta r)$ for RQ2, the relative cost of the atomic variant in
the browser over the same relative cost natively.
\item[D7, exclusions.] No statistical outlier removal. Only three exclusions
exist, decided in advance and counted: failed runs, runs that fail strict
validation, and runs whose foreign CPU time is above 0.25~s.
\item[D8, multiple comparisons.] The \emph{primary family} has ten
contrasts, all at $N=8$: $\Delta r$ for $k \in \{1,4,16,256\}$ and both
browsers (eight), and the filter's log speedup ratio for both browsers
(two). Each has a two-sided bootstrap $p$-value, twice the smaller share of
bootstrap estimates on one side of zero, floored at $1/10\,000$. The
decision is Holm's procedure at a family-wise $\alpha=0.05$. Primary
contrasts are shown with Bonferroni 99.5\% intervals, which hold
simultaneously for the family. Everything else is exploratory: uncorrected
95\% intervals, labelled as such.
\end{description}

\paragraph{Why these choices.} Intervals on ratios show the size of an
effect and its uncertainty together, which a bare test does not. The
bootstrap makes no assumption about the shape of the run times. The
geometric mean is the natural centre for quantities compared by ratio. A
small primary family, fixed in advance, keeps the correction for multiple
comparisons from being diluted by the hundreds of descriptive cells.

\paragraph{Amendments.} Two were made before any data: every scaling
comparison is read next to the absolute time ratio
$T(\mathrm{browser},N)/T(\mathrm{native},N)$, because a relative speedup can
look better only because its baseline is slow~\cite{hoefler2015}; and a
configuration whose run value is under 100 clock ticks is flagged. One was
made after the first campaign's primary results existed, and is labelled as
such: a \emph{robustness check} that tests each primary contrast with
Mann-Whitney $U$~\cite{mann1947} and Cliff's $\delta$ on per-run ratios
($\log(T_1/T_N)$ for RQ1, $\log(T_{\mathrm{atomic}}/T_{\mathrm{local}})$ for
RQ2), Holm-adjusted. The two runs of a ratio are separate processes, paired
by repetition index, so the check is repeated over 1000 random re-pairings.
The bootstrap decisions stay the reported ones; a disagreement would be
reported as not robust.

\paragraph{Sensitivity checks.} The plan names three: the per-run mean
instead of the median, the median across runs instead of the geometric
mean, and no exclusion by foreign CPU time. A conclusion that changes under
one of them is reported as not robust.

\paragraph{The data check.} Before any analysis a report counts the runs
done, missing and failed, the checksum conflicts, the runs failing strict
validation, the runs above the threshold, the configurations under the
clock factor, and the configurations whose run values trend with the
execution order (Spearman's $|\rho| > 0.5$ between position and value; with
30 runs and no drift this has a probability of about 0.5\%, so about one
configuration of the main campaign can be flagged by chance). What is done
with each finding was fixed beforehand: a conflict or a missing run stops
the analysis; a trending configuration is reported and kept, since the
random order turns drift into variance, not bias; no outlier is removed.
Every analysis command reads a campaign through one loader that applies
these rules.

\paragraph{Post hoc tables.} Three descriptive tables were added after the
plan and decide nothing: the two factors of each RQ2 ratio
(Section~\ref{sec:decomposition}), throughput with the absolute speedup,
and the Karp-Flatt fraction.

\paragraph{The explanations have their own rule.} For the questions of
Section~\ref{sec:why}, each hypothesis, the test of it and what the test
would show if the hypothesis held or failed were written in a dated record
before the test ran. Without that, any outcome can be read as support.
